\documentclass[10pt,a4paper]{amsart}
\usepackage[margin=19mm]{geometry}
\usepackage{amsmath,amssymb,mathtools}
\usepackage[scr=rsfs,cal=euler]{mathalfa}
\usepackage{xfrac}
\usepackage[unicode,hidelinks,bookmarks=false]{hyperref}
\newcommand{\ie}{i.e.\ }
\newcommand{\cf}{cf.\ }
\newcommand{\resp}{respectively\ }
\newcommand{\Subsection}{\subsection}
\providecommand{\nobreakdash}{\nobreak\hbox{-}}
\setlength{\emergencystretch}{2em}
\multlinegap0pt
\usepackage{bm}
\usepackage{bbm}
\usepackage{dsfont}
\usepackage{xspace}
\usepackage{cancel}
\usepackage{mathtools}
\usepackage{amsthm}
\makeatletter
\@ifundefined{thm}{\newtheorem{thm}{Thm}}{}
\@ifundefined{prop}{\newtheorem{prop}[thm]{Proposition}}{}
\makeatother
\title[Enforcing TSP-Optimality in Fair Vehicle Routing by Cutting ]{Enforcing TSP-Optimality in Fair Vehicle Routing by Cutting Planes}
\author{Bart van Rossum, Rui Chen, Andrea Lodi}
\date{Source: arXiv:2604.23748v1 (CC BY 4.0)}
\begin{document}
\maketitle

\noindent\emph{Source and licence.} Enforcing TSP-Optimality in Fair Vehicle Routing by Cutting Planes. Version arXiv:2604.23748v1 (CC BY 4.0), distributed under the Creative Commons Attribution 4.0 International licence, as recorded in the arXiv licence metadata for this exact version and shown on the abstract page https://arxiv.org/abs/2604.23748v1. Full rights notice: licenses/arxiv-2604.23748-CC-BY-4.0.txt.\par\smallskip

\noindent\textit{Editorial scope.} Counted unit: the Prop whose source label is \texttt{None} in the article (source0.tex lines 295--297, source label \texttt{None}), delivered together with the article's own complete original proof of it (source0.tex lines 298--315, one \texttt{proof} environment of 18 lines). No mathematics is written, repaired or re-derived: statement and proof are the article's own text line for line. Required context retained uncounted: aspt:depot (lines 254--256); eq:lifted\_cut (lines 291--293). Every definition, display and label of those lines is retained unchanged. Omitted from this package and not redistributed: 1--253, 257--290, 294--294, 316--366. Nothing inside a retained excerpt is omitted or altered: every retained body is delivered exactly as the article's lines are written, so no omission receipt of any kind accompanies this package. Citations: the retained excerpts carry no citation command at all, so no bibliography entry of the article is transcribed and no reference is redistributed. Reference adaptation: none was needed and none was made; every reference inside the delivered unit resolves inside it, so no unresolved reference or citation is produced. Printed slips kept literal: no printed word, symbol, hypothesis or number of the counted statement or of its proof was altered. Layout and class: the article's own class is \texttt{cas-dc} with packages including natbib, tikz, adjustbox, geometry, todonotes, amsmath, amsthm, bm, pgfplots, subfiles, subcaption; the wrapper is the lane's pinned \texttt{amsart} editorial wrapper at 10pt on A4, which restates the theorem-like environments the retained text needs and adds the article's own macro definitions that the retained text needs (none), with extra wrapper packages (bm, bbm, dsfont, xspace, cancel, mathtools). The delivered numbering is the unit's own. Assets: no retained span contains a figure environment, a graphics-inclusion command, a PDF-page inclusion, a table or a TikZ picture; no image file is copied into this package, the package declares no asset dependency and no image file is redistributed with it. Licence: version arXiv:2604.23748v1 of this article is distributed under the Creative Commons Attribution 4.0 International licence, as recorded in the arXiv licence metadata for this exact version and shown on the abstract page https://arxiv.org/abs/2604.23748v1. A full rights notice is delivered at licenses/arxiv-2604.23748-CC-BY-4.0.txt. Characters: every retained excerpt is pure ASCII, so the delivered engine, XeLaTeX with its default Latin Modern text font, reports no missing character.
\begin{equation}
    v_2, \ldots, v_{p-1} \neq 0. \label{aspt:depot}
\end{equation}

\begin{equation}
x(A(P))+ x\!\big(\bigcup_{h=1}^{p-1} \Delta^{+}_{h}(P)\big)\leq |A(P)|-1. \label{eq:lifted_cut}
\end{equation}

\begin{prop}
For a given TSP-violating path $P$, the forward-lifted cut~\eqref{eq:lifted_cut} is valid for F-CVRP-TSP.
\end{prop}

\begin{proof}
Assume to the contrary that there exists a feasible solution $x$ violating~\eqref{eq:lifted_cut}, i.e., for which the left-hand side equals at least $|A(P)|$. We prove by induction that $x_{v_h, v_{h+1}} = 1$ for $h = 1, \ldots, p-1$, which implies $x(A(P)) = |A(P)|$, contradicting TSP-optimality of $x$.

Observe that the cut can be rewritten as
\begin{equation}
\sum_{h=1}^{p-1} \big( x_{v_h, v_{h+1}} + x(\Delta^{+}_{h}(P)) \big) \leq |A(P)| - 1. \label{eq:lifted_cut_rewritten}
\end{equation}
%We note two facts. First, $\Delta^{+}_{1}(P) = \emptyset$, since the partial path $(v_1)$ consists of a single node and no lifting arcs can be defined. Second,
We note three facts. First, by \eqref{aspt:depot}, the node $v_h$ is a customer node for all $h = 2, \ldots, p-1$. By the degree constraint, it follows that $x_{v_h, v_{h+1}} + x(\Delta^{+}_{h}(P)) \leq 1$ for all such $h$. Second, $\Delta^{E,+}_{1}(P) = \Delta^{TSP,+}_{1}(P) = \emptyset$, since the partial path $(v_1)$ consists of a single node. Third, $x_a=0$ for $a\in \Delta^{Q,+}_{1}(P)$ by feasibility of $x$ and the definition of $\Delta^{Q,+}_{1}(P)$.

Hence, each of the $p-1$ terms in~\eqref{eq:lifted_cut_rewritten} contributes at most $1$, and the first term contributes at most $x_{v_1, v_2}$. For the left-hand side to reach $|A(P)| = p - 1$, it must hold that $x_{v_1, v_2} = 1$ and $x_{v_h, v_{h+1}} + x(\Delta^{+}_{h}(P)) = 1$ for all $h = 2, \ldots, p-1$.

We now apply induction on $h$. The base case $h=1$ requires $x_{v_1, v_2} = 1$, as already argued above.

For the induction step, suppose $x_{v_h, v_{h+1}} = 1$ for all $h = 1, \ldots, k-1$ with $2 \leq k \leq p-1$. We aim to show that $x_{v_k, v_{k+1}} = 1$. As the cut is violated, we have $x_{v_k, v_{k+1}} + x(\Delta^{+}_{k}(P)) = 1$. Suppose $x_{v_k, v_{k+1}} = 0$. Then $x(\Delta^{+}_{k}(P)) = 1$, meaning there exists an arc $(v_k, j) \in \Delta^{+}_{k}(P)$ with $x_{v_k, j} = 1$. By the induction hypothesis, the partial path $(v_1, \ldots, v_k)$ is traversed by $x$. However, by construction of $\Delta^{+}_{k}(P)$, extending this path with arc $(v_k, j)$ violates elementarity, TSP-optimality, or the capacity constraint. In each case, this contradicts feasibility of $x$. Hence, $x_{v_k, v_{k+1}} = 1$.

By induction, we conclude that path $P$ is fully traversed, contradicting TSP-optimality of $x$.
\end{proof}

\end{document}
